\chapter{Einführung}
\label{c:intro}

\section{Motivation}
\label{s:intro:mot}

Die Techdivision ist intern in sogenannten Services strukturiert, die häufig ein intensives Zusammenarbeiten erfordern. Diese Zusammenarbeit findet über verschiedene Kommunikationskanäle statt, darunter JIRA Boards und Support-Chats. Oft werden Routinefragen gestellt, die automatisiert werden könnten. 

Durch den Aufstieg von verschiedenen textbasierten Chat-Modellen wie OpenAIs Chat-GPT oder Googles Gemini, hat man in den letzten Jahren gesehen, zu was Modelle inzwischen in der Lage sind. Sie sind nicht mehr nur in der Lage, mit Menschen zu kommunizieren, sondern helfen diesen auch bei ihren Jobs oder Schulangelegenheiten. Verschiedene trainierte Modelle sind inzwischen in der Lage, komplizierte Probleme zu lösen und Menschen das Leben einfacher zu machen. 

Dadurch stellt sich speziell in den Support-Chats die Frage, wie viele der Supportanfragen sich durch einen KI-gesteuerten Chatbot ohne Hilfe eines Mitarbeiters lösen lassen würden und inwiefern ein Chatbot zur Entlastung der Teams und zur schnelleren Lösung von Problemen beitragen könnte.
Durch die Automatisierung einfacher Supportanfragen könnte zum einen Zeit gespart, zum Anderen auch die Ressourcen der Teams effektiver genutzt werden. Mitarbeiter könnten sich verstärkt auf die schwierigeren Supportanfragen konzentrieren, während der Chatbot die repetitiven und einfachen klärt, was nicht nur das Team entlastet, sondern auch die Reaktionszeit auf Supportanfragen erhöht.

\section{Zielsetzung}
\label{s:intro:goal}
Das Ziel der Arbeit besteht darin, einen Chatbot zu konzipieren und einen Prototyp zu entwickeln, der mit Hilfe der Daten der internen Datenquellen der Techdivision Support leistet und Routinefragen beantwortet. Der Chatbot soll idealerweise innerhalb von Google Chat integriert werden. Dabei ist es wichtig, dass der Chatbot, statt in Einzelchats, in einer Chat Gruppe läuft. Diese beinhaltet neben den Support Suchenden, ebenfalls Mitarbeiter des jeweiligen Services. Auf diese Weise kann, falls notwendig, extern Eingegriffen werden. Außerdem sollte der Chatbot in der Lage sein zu überprüfen, wie viele der aktuellen Supportanfragen er eigenständig beantworten oder lösen kann. Falls der Chatbot keine Antwort auf eine Anfrage haben sollte, soll er in den angebundenen Datenquellen nachschlagen können, um mit dem Wissen seine Wissensbasis zu erweitern und die Fragen entsprechend zu beantworten.